\begin{itemize}
    \item \textbf{Study population: }
% Describe key characteristics of the target study population, e.g. adults 18-50 years of age who are HIV-infected

    \item \textbf{Treatment Allocation and Randomization: }
% stratified randomization, minimization, non-randomized concurrent control, or single arm 

    \begin{itemize}
        \item \textbf{allocation ratio}  %  = 1:1, 2:1... 
        \item \textbf{stratification factors: } %none, disease stage(I, II, and III), study site, region
    \end{itemize}

    \item \textbf{Blinding: } % open label, single blind, double blind
% If blinded, be specific about who has access to treatment allocation information 
% prior to database lock and who is/are blinded, 
% e.g., subject, investigator, surgeon, or outcome evaluator). 
% Delete this bullet point if single arm.

    \item \textbf{Treatment Arms: }
% Delete this bullet point if single arm; 
% add lines below for >2 arms; can use a table; 
% include regimen information (dose, route, schedule) if necessary

    \begin{itemize}
        \item \textbf{experimental: } % e.g. gene therapy
        \item \textbf{control: } % e.g. standard of care, placebo
    \end{itemize}

    \item \textbf{Site(s): } 
% number of sites ≈ x in US, multi-regions (US, UK, Japan, ...). If available, note US study population percentage

    \item \textbf{Other features: }
% Describe unusual and important design features not listed above, 
% e.g., subjects will not receive product till two weeks after randomization and product manufacturing, 
% or lot-consistency, as a sub-study in the overall study design, 
% needs to be demonstrated before combining lots for efficacy comparison with the control group; 
% add an overall study schematic flow chart here if needed for complicated study design. 
% Delete this bullet point if none. DO NOT describe interim looks/non-inferiority design in this section.

\end{itemize}